% Source: https://artofproblemsolving.com/community/c3h76230_2000_national_team_round
% Stable identity: aops-topic-76230-q01
Rotating a right triangle about a leg gives a cone whose height is that leg and whose base radius is the other leg.

\begin{asy}
usepackage("amsmath");
size(8cm);
defaultpen(fontsize(10pt));
path ra(pair P,pair u,pair v,real s=0.2){return P+s*unit(u)--P+s*unit(u)+s*unit(v)--P+s*unit(v);}
pair A=(0,0),B=(4,0),T=(0,3);
draw((0,-1.1)--(0,3.9),gray+dashdotted); label("axis",(0,3.9),N,fontsize(8pt));
draw(A--B--T--cycle,linewidth(0.9)); draw(A--(-4,0)--T,gray+dashed);
draw(ellipse((0,0),4,0.6),blue);
draw(ra(A,B-A,T-A,0.35));
label("$3$",(0,1.5),W); label("$4$",(2,-0.6),S); label("$5$",(2,1.5),NE);
pair o=(10,0);
pair A2=o,B2=o+(3,0),T2=o+(0,4);
draw(o+(0,-1.1)--o+(0,4.9),gray+dashdotted); label("axis",o+(0,4.9),N,fontsize(8pt));
draw(A2--B2--T2--cycle,linewidth(0.9)); draw(A2--o+(-3,0)--T2,gray+dashed);
draw(ellipse(o,3,0.5),blue);
draw(ra(A2,B2-A2,T2-A2,0.35));
label("$4$",o+(0,2),W); label("$3$",o+(1.5,-0.5),S); label("$5$",o+(1.5,2),NE);
label("\mbox{}",(min(currentpicture,user=true)+max(currentpicture,user=true))/2,black); // keeps the text after the figure black when LaTeX typesets the labels inline
\end{asy}

About the leg of length $3$: radius $4$, height $3$, volume $\dfrac{1}{3}\pi\cdot4^2\cdot3=16\pi\approx50.27$. About the leg of length $4$: radius $3$, height $4$, volume $\dfrac{1}{3}\pi\cdot3^2\cdot4=12\pi\approx37.70$. The greater volume is $16\pi\approx\boxed{50}$ cubic centimetres. (The answer $30$ suggested in the post is wrong.)
